\documentclass[12pt, twoside]{article}
\input{/Users/chris/github/course-files/docs/ib/preamble}

\fancyhead[LE]{\thepage}
\fancyhead[RO]{Markscheme}
\fancyhead[LO]{La Scuola d'Italia / Dr.\ Huson / IB Math 12 HL \\* 10 November 2026}

\begin{document}
\subsubsection*{5.5 Classwork: Proof, Differential Equations and Series --- Markscheme}

Total 59 marks.

\begin{enumerate}[itemsep=0.3cm]

%--- Problem 1 [13 marks] ---
\item[1.] [Maximum mark: 13]

\textbf{(a)} $n^{3} - n = (n - 1)n(n + 1)$ \hfill \textbf{A1}

$n - 1$ and $n$ are consecutive integers, so one of them is even, and
the product is even \hfill \textbf{R1}

\emph{Accept a case argument: $n = 2m$ gives $2m(4m^{2} - 1)$;
$n = 2m + 1$ gives $(2m + 1)(2m)(2m + 2)$. Checking a few values of $n$
is not a proof: A0R0.}

\textbf{(b)} Assume there exists $\alpha \in \mathbb{Z}$ such that
$\alpha^{3} + \alpha + 1 = 0$ \hfill \textbf{M1}

\emph{M0 for vague openings such as ``consider the equation has integer
roots'' or ``let $\alpha \in \mathbb{Z}$ be a root''; the assumption must
state that an integer root exists. Condone $x$ in place of $\alpha$.
Later marks do not depend on this M1 (as in the IB markscheme for 2022
May P1 TZ2 Q9).}

$\alpha^{3} + \alpha = -1$ \hfill \textbf{M1}

$\alpha^{3} + \alpha = \left(\alpha^{3} - \alpha\right) + 2\alpha$
\hfill \textbf{A1}

$\alpha^{3} - \alpha$ is even by part (a) and $2\alpha$ is even, so
$\alpha^{3} + \alpha$ is even \hfill \textbf{R1}

but $-1$ is odd, a contradiction, so the equation has no integer roots
\hfill \textbf{R1}\,\textbf{AG}

\emph{Accept $\alpha^{3} + \alpha = \alpha\left(\alpha^{2} + 1\right)$ with
the case argument ($\alpha$ even, or $\alpha$ odd so $\alpha^{2} + 1$
even) for A1R1. Accept the method of the increasing function:
$f'(x) = 3x^{2} + 1 > 0$, so $f$ is increasing; $f(-1) = -1$ and
$f(0) = 1$, so the only real root lies between $-1$ and $0$, which
contradicts the assumption: M1 (assumption) M1A1 (derivative and
increasing) R1 (the two values) R1 (contradiction).}

\textbf{(c)} Assume there exist $p, q \in \mathbb{Z}$ such that
$p^{2} - 8q - 3 = 0$ \hfill \textbf{M1}

$p^{2} = 8q + 3 = 2(4q + 1) + 1$, which is odd \hfill \textbf{A1}

so $p$ is odd (if $p$ were even, $p^{2}$ would be even) \hfill \textbf{R1}

let $p = 2k + 1$, $k \in \mathbb{Z}$: \
$4k^{2} + 4k + 1 = 8q + 3$ \hfill \textbf{M1}

$4k^{2} + 4k - 8q = 2$

$2\left(k^{2} + k - 2q\right) = 1$ \hfill \textbf{A1}

the left side is even and $1$ is odd, a contradiction, so
$p^{2} - 8q - 3 \neq 0$ \hfill \textbf{R1}\,\textbf{AG}

\emph{Accept $4k(k + 1) = 8q + 2$: by part (a) reasoning $k(k + 1)$ is
even, so the left side is a multiple of $8$, but $8q + 2$ is not. Accept
$k^{2} + k - 2q = \frac{1}{2}$, not an integer.}

\newpage

%--- Problem 2 [17 marks] ---
\item[2.] [Maximum mark: 17]

\textbf{(a)} $\displaystyle\int \left(x^{2} + 1\right)e^{x}\,\mathrm{d}x
= \left(x^{2} + 1\right)e^{x} - \int 2xe^{x}\,\mathrm{d}x$ \hfill \textbf{M1A1}

$\displaystyle\int 2xe^{x}\,\mathrm{d}x = 2xe^{x} - \int 2e^{x}\,\mathrm{d}x
= 2xe^{x} - 2e^{x}$ \hfill \textbf{M1}

$\left(x^{2} - 2x + 3\right)e^{x} + C$ \hfill \textbf{A1}

\emph{Accept an unsimplified equivalent. Condone a missing $+C$ here; it
is needed in (b).}

\textbf{(b)} $\dfrac{\mathrm{d}y}{\mathrm{d}x} + y = x^{2} + 1$;
integrating factor $e^{\int 1\,\mathrm{d}x} = e^{x}$ \hfill \textbf{M1}

$\dfrac{\mathrm{d}}{\mathrm{d}x}\left(ye^{x}\right) = \left(x^{2} + 1\right)e^{x}$
\hfill \textbf{A1}

$ye^{x} = \left(x^{2} - 2x + 3\right)e^{x} + C$ \hfill \textbf{A1}

$x = 0,\ y = 1$: $1 = 3 + C$, so $C = -2$ \hfill \textbf{M1}

$y = x^{2} - 2x + 3 - 2e^{-x}$ \hfill \textbf{A1}\,\textbf{AG}

\emph{The final A1 needs the division by $e^{x}$ applied to the constant:
$Ce^{-x}$. Common error: $y = x^{2} - 2x + 3 + C$ (constant not divided
by $e^{x}$) leads to $y = x^{2} - 2x + 1$, so the AG is not reached:
A0 for the last two marks unless the error is corrected. FT the
integral from (a).}

\textbf{(c)} $y'(0) = 0 + 1 - 1 = 0$ \hfill \textbf{A1}

$y'' = 2x - y'$ \hfill \textbf{M1}

$y''(0) = 0$ \hfill \textbf{A1}

$y''' = 2 - y''$, so $y'''(0) = 2$ \hfill \textbf{A1}

$y = 1 + \dfrac{2}{3!}x^{3} = 1 + \dfrac{x^{3}}{3}$ \hfill \textbf{A1}

\emph{Accept $1 + 0x + 0x^{2} + \frac{1}{3}x^{3}$. Expanding the solution
from (b) instead is not what the question asks: M0, but award the final
A1 for $1 + \frac{x^{3}}{3}$.}

\textbf{(d)} $y' = 2x - 2 + 2e^{-x}$, \ $y'' = 2 - 2e^{-x}$ \hfill \textbf{M1}

$y''(0) = 2 - 2 = 0$ \hfill \textbf{A1}

$y'' < 0$ for $x < 0$ (since $e^{-x} > 1$) and $y'' > 0$ for $x > 0$, so
the concavity changes at $x = 0$; and $y(0) = 3 - 2 = 1$ \hfill \textbf{R1}

\emph{R1 needs the sign change of $y''$, not only $y''(0) = 0$. Accept
the series from (c): near $0$, $y \approx 1 + \frac{x^{3}}{3}$, so
$y'' \approx 2x$, which changes sign at $0$.}

\newpage

%--- Problem 3 [15 marks] ---
\item[3.] [Maximum mark: 15]

\textbf{(a)} $n = 1$: $f'(x) = \dfrac{1}{1 + x}$ and
$\dfrac{(-1)^{0}\,0!}{(1 + x)^{1}} = \dfrac{1}{1 + x}$, so true for
$n = 1$ \hfill \textbf{A1}

assume true for $n = k$:
$f^{(k)}(x) = \dfrac{(-1)^{k - 1}(k - 1)!}{(1 + x)^{k}}$ \hfill \textbf{M1}

\emph{Do not accept ``let $n = k$'' or ``assume $n = k$''; the
statement must be assumed true.}

$f^{(k + 1)}(x) = \dfrac{\mathrm{d}}{\mathrm{d}x}
\left[(-1)^{k - 1}(k - 1)!\,(1 + x)^{-k}\right]$ \hfill \textbf{M1}

$= (-1)^{k - 1}(k - 1)!\,(-k)(1 + x)^{-k - 1}$ \hfill \textbf{A1}

$= \dfrac{(-1)^{k}\,k!}{(1 + x)^{k + 1}}$ \hfill \textbf{A1}

this is the statement for $n = k + 1$ \hfill \textbf{R1}

since true for $n = 1$, and true for $n = k$ implies true for
$n = k + 1$, the statement is true for all $n \in \mathbb{Z}^{+}$
\hfill \textbf{R1}

\emph{The final R1 needs both the base case and the inductive step to
have been attempted.}

\textbf{(b)} $f(0) = 0$ and $f^{(n)}(0) = (-1)^{n - 1}(n - 1)!$
\hfill \textbf{M1}

coefficient of $x^{n}$: $\dfrac{(-1)^{n - 1}(n - 1)!}{n!}$ \hfill \textbf{A1}

$= \dfrac{(-1)^{n - 1}}{n}$, so
$f(x) = x - \frac{x^{2}}{2} + \frac{x^{3}}{3} - \cdots
+ \frac{(-1)^{n - 1}x^{n}}{n} + \cdots$ \hfill \textbf{A1}\,\textbf{AG}

\textbf{(c)} $\ln\left(\dfrac{1 + x}{1 - x}\right) = \ln(1 + x) - \ln(1 - x)$
\hfill \textbf{M1}

$\ln(1 - x) = -x - \frac{x^{2}}{2} - \frac{x^{3}}{3} - \frac{x^{4}}{4}
- \frac{x^{5}}{5} - \cdots$ \hfill \textbf{A1}

$2x + \dfrac{2x^{3}}{3} + \dfrac{2x^{5}}{5}$ \hfill \textbf{A1}

\emph{Common error: $\ln\left(\frac{1 + x}{1 - x}\right)$ written as
$\frac{\ln(1 + x)}{\ln(1 - x)}$, or $\ln(1 + x)$ split as
$\ln 1 + \ln x$: M0.}

\textbf{(d)} $\dfrac{\frac{2x^{3}}{3} + \frac{2x^{5}}{5} + \cdots}{x^{3}}
= \dfrac{2}{3} + \dfrac{2x^{2}}{5} + \cdots$ \hfill \textbf{M1}

limit $= \dfrac{2}{3}$ \hfill \textbf{A1}

\emph{FT from (c). Accept l'H\^opital's rule applied correctly.}

\newpage

%--- Problem 4 [14 marks] ---
\item[4.] [Maximum mark: 14]

\textbf{(a)} $w + w^{*} = (a + bi) + (a - bi) = 2a = 2\,\mathrm{Re}(w)$
\hfill \textbf{A1}

$ww^{*} = (a + bi)(a - bi) = a^{2} - b^{2}i^{2}$ \hfill \textbf{M1}

$= a^{2} + b^{2} = |w|^{2}$ \hfill \textbf{A1}\,\textbf{AG}

\textbf{(b)} $w^{*} = re^{-i\theta}$ \hfill \textbf{A1}

$\left(w^{*}\right)^{n} = r^{n}e^{-in\theta}$ \hfill \textbf{M1}

$w^{n} = r^{n}e^{in\theta}$, so
$\left(w^{n}\right)^{*} = r^{n}e^{-in\theta} = \left(w^{*}\right)^{n}$
\hfill \textbf{A1}\,\textbf{AG}

\emph{Accept the same argument in cis form:
$\left(r\,\mathrm{cis}(-\theta)\right)^{n} = r^{n}\,\mathrm{cis}(-n\theta)$.}

\textbf{(c)} $1 - i = (1 + i)^{*}$, so $(1 - i)^{n} = \left((1 + i)^{n}\right)^{*}$
by (b) \hfill \textbf{R1}

$(1 + i)^{n} + (1 - i)^{n} = z + z^{*} = 2\,\mathrm{Re}(z)$, where
$z = (1 + i)^{n}$, which is real by (a) \hfill \textbf{R1}\,\textbf{AG}

\textbf{(d)} $1 + i = \sqrt{2}\,e^{i\frac{\pi}{4}}$ \hfill \textbf{A1}

$(1 + i)^{n} = 2^{\frac{n}{2}}e^{i\frac{n\pi}{4}}$ \hfill \textbf{M1}

$(1 + i)^{n} + (1 - i)^{n} = 2\,\mathrm{Re}\left(2^{\frac{n}{2}}e^{i\frac{n\pi}{4}}\right)$
\hfill \textbf{M1}

$= 2 \times 2^{\frac{n}{2}}\cos\dfrac{n\pi}{4} = 2^{\frac{n}{2} + 1}\cos\dfrac{n\pi}{4}$
\hfill \textbf{A1}\,\textbf{AG}

\emph{Accept adding $2^{\frac{n}{2}}e^{i\frac{n\pi}{4}}$ and
$2^{\frac{n}{2}}e^{-i\frac{n\pi}{4}}$ directly, the sine terms cancelling,
for the second M1.}

\textbf{(e)} $2^{7}\cos 3\pi$ \hfill \textbf{M1}

$-128$ \hfill \textbf{A1}

\emph{Accept $(1 + i)^{2} = 2i$, so $(1 + i)^{12} = (2i)^{6} = -64$, and the
same for $(1 - i)^{12}$, giving $-128$.}

\end{enumerate}
\end{document}
